% FILE: 02_lineage_and_context.tex
% Project: sources-wasm4pm-compat
% Role: type-foundry-and-witness-lattice-authority

\section{Lineage and Context}
\label{sec:sources-wasm4pm-compat-lineage}

\subsection{2016 Language-Model Lesson}
\label{sec:sources-wasm4pm-compat-lineage-2016}

The 2016 language-model experiment established a foundational negative result for this
research program: local text imitation does not conserve consequence. A model that
produces text indistinguishable from expert output can simultaneously fail to preserve
the causal structure that made the expert output meaningful. In the process intelligence
context, this translates directly: a system that emits conformance scores, fitness values,
or alignment costs without binding those values to a receipted, non-forgeable evidence
chain is performing imitation, not measurement.

The type-law foundry is the architectural response to this lesson. The
\texttt{Evidence<T, State, Witness>} container is designed so that a conformance score
cannot be emitted without simultaneously being bound to the event-log payload
that produced it (type parameter \texttt{T}), the process-model state at the point of
measurement (\texttt{State}), and the lattice-monotonic proof of how that state was reached
(\texttt{Witness}). The BLAKE3 hash over all three fields, plus the epoch nonce and
Ed25519 signature, makes it cryptographically impossible to separate the score from its
evidence chain post-construction. The 2016 lesson---that imitation without consequence
preservation is epistemologically void---is encoded as a compile-time type law.

\subsection{Enterprise Process Gap}
\label{sec:sources-wasm4pm-compat-lineage-enterprise}

Enterprise process intelligence systems face a structural gap between the process models
they declare and the processes they actually execute. This gap is not primarily a tooling
problem; it is an evidence problem. When an organization claims that its order-to-cash
process has a fitness of 0.94 against a BPMN model, that claim is only as strong as the
provenance chain connecting the claim to the raw event log, the process model, the
conformance checking algorithm, and the authority that signed the result.

Without a cryptographic binding between those components, the claim is a narrative, not
a receipt. Under M\&A due diligence conditions---where the claim may be worth millions
in synergy projections---a narrative is not board-admissible. The \texttt{wasm4pm-compat}
type-law foundry was designed specifically to make board-admissible claims manufacturable:
claims that carry a BLAKE3 hash, an Ed25519 signature from an auditor role, a fitness
value proven to be in $[0, 1]$ via the \texttt{Between01<N,D>} const-generic rational
bound, and a \texttt{WitnessMarker} at or above the threshold for the claimed process
family.

\subsection{Relationship to the Chatman Equation}
\label{sec:sources-wasm4pm-compat-lineage-chatman}

The Chatman Equation $A = \mu(O^*)$ asserts that autonomic actuation $A$ is a function
of the measurement operator $\mu$ applied to an observed process trace $O^*$.
\texttt{wasm4pm-compat} provides the type-law infrastructure that makes $\mu$ computable
in a non-forgeable way. The \texttt{WitnessMarker::BlueRiverDam} lattice position---the
top element \texttt{\(\top\)} of the witness lattice, carrying an \texttt{AutonomicMarking}
with MAPE-K enforcement---is the typed representation of the moment when $\mu(O^*)$
has been computed, receipted, and is ready to drive actuation.

Critically, the lattice structure ensures that \texttt{BlueRiverDam} cannot be reached
without passing through the lower lattice positions. A caller cannot mint a
\texttt{WitnessMarker::BlueRiverDam} evidence block without first establishing that the
trace has been parsed (bottom), validated and replayed (middle), and signed by an auditor
role. The join-to-top condition---where \texttt{witness\_join} reaches the top element and
triggers execution halt---is the type-level expression of the equation's boundary condition:
when the measurement process has consumed all available evidence, autonomous action begins.

\subsection{Relationship to Receipts and Replay}
\label{sec:sources-wasm4pm-compat-lineage-receipts}

The receipt machinery in \texttt{wasm4pm-compat} is the operational embodiment of the
research program's claim that every authoritative assertion must be receipted. The
\texttt{ArtifactReceipt} struct in \texttt{compat/manufacturing/receipt\_ledger.rs}
binds a BLAKE3 content hash, a \texttt{WitnessMarker}, a lifecycle state (always
\texttt{Sealed} at time of recording), a Unix timestamp, and a compliance audit result.
The \texttt{ReceiptLedger} maintains a \texttt{HashMap}-backed provenance registry with
\texttt{verify\_all()} and \texttt{count\_by\_witness()} methods.

The manufacturing manifest \texttt{manufacturing-manifest-20260601.yaml} records that
three receipts were issued on 2026-06-01: \texttt{witness-markers-20260601-receipt}
(witness: BlueRiverDam), \texttt{process-forms-20260601-receipt} (witness: VanDerAalst1998),
and \texttt{boundary-law-20260601-receipt} (witness: BlueRiverDam). All three carry
Ed25519 signature status ISSUED. The concept of a receipt-shaped object---an
\texttt{Evidence} block where \texttt{State} is terminal, \texttt{Witness} carries
fitness $\geq$ threshold, and the hash is signed by an auditor role---directly connects
the type-law architecture to the downstream replay and conformance checking operations
that produce the fitness values those receipts attest to.
